\begingroup
\fontsize{8.5pt}{9.5pt}\selectfont
\renewcommand{\arraystretch}{1.3}
\setlength{\tabcolsep}{4pt}

\begin{longtable}{|
	>{\centering\arraybackslash}p{1.1cm}|
	>{\raggedright\arraybackslash}p{3.2cm}|
	>{\raggedright\arraybackslash}p{9.2cm}|}

	\caption{Requisitos definidos para el servicio de directorio}
	\label{tab:requisitos}
	\\
	\hline
	\tableheader
	\textbf{ID}                         &
	\textbf{Requisito}                  &
	\textbf{Definición}
	\\
	\hline
	\endfirsthead

	\hline
	\tableheader
	\textbf{ID}                         &
	\textbf{Requisito}                  &
	\textbf{Definición}
	\\
	\hline
	\endhead

	R-01                                &
	Centralización de identidades       &
	El servicio deberá permitir centralizar la gestión de las identidades de los usuarios y facilitar su utilización en los diferentes servicios tecnológicos, manteniendo mecanismos de autorización diferenciados para cada sistema.
	\\
	\hline

	R-02                                &
	Gestión del ciclo de vida           &
	El servicio deberá permitir gestionar el ciclo de vida de las cuentas de usuario, incluyendo actividades de creación, modificación, verificación, revocación y eliminación o expiración de las credenciales.
	\\
	\hline

	R-03                                &
	Políticas de seguridad              &
	El servicio deberá permitir definir y aplicar políticas de seguridad para la gestión de usuarios y credenciales, de acuerdo con las reglas establecidas para el entorno.
	\\
	\hline

	R-04                                &
	Control y monitoreo de accesos      &
	El servicio deberá proporcionar mecanismos para controlar y monitorear los accesos de los usuarios y facilitar la revisión de los permisos asignados.
	\\
	\hline

	R-05                                &
	Auditoría y trazabilidad            &
	El servicio deberá registrar las actividades relevantes relacionadas con la gestión identidades y el control de accesos, de manera que puedan ser consultadas para fines de auditoría.
	\\
	\hline

	R-06                                &
	Fortalecimiento de la autenticación &
	El servicio deberá proporcionar mecanismos para fortalecer la autenticación de los usuarios mediante políticas de credenciales y mecanismos adicionales de autenticación, de acuerdo con las necesidades de cada servicio.
	\\
	\hline

	R-07                                &
	Recuperación de credenciales        &
	El servicio deberá proporcionar mecanismos seguros para la recuperación o restablecimiento de las credenciales de los usuarios.
	\\
	\hline

	R-08                                &
	Evolución y ampliación              &
	El servicio deberá permitir su evolución y ampliación ante la incorporación futura de nuevos sistemas o activos que requieran mecanismos de autenticación y control de acceso.
	\\
	\hline
\end{longtable}
\endgroup
